
\documentclass[11pt]{article}

\usepackage{amsmath,amssymb,amsthm,geometry}
\usepackage{hyperref}
\geometry{margin=1in}

\title{Radiation, Constraints, and the Non-Existence of Isolated Systems in Relativistic Field Theory}
\author{}
\date{}

\begin{document}

\maketitle

\begin{abstract}
Although Maxwell’s equations admit source-free wave solutions, radiation is not an autonomous physical phenomenon. This paper argues that electromagnetic radiation arises only through global boundary conditions, interaction with matter, and the asymptotic structure of spacetime. By synthesizing results from constrained Hamiltonian dynamics, gauge theory, the no-interaction theorem, Haag’s theorem, absorber theory, and gravitational radiation theory, we show that radiation is fundamentally relational and asymptotic rather than local and self-sustaining. This perspective clarifies the physical meaning of free fields, the role of time asymmetry, and the impossibility of isolated subsystems in relativistic physics.
\end{abstract}

\section{Introduction}

Maxwell’s equations admit nontrivial solutions in the absence of sources, often interpreted as electromagnetic waves propagating freely through space. This mathematical fact has led to the widespread belief that radiation is an autonomous physical entity. In this paper we argue that this interpretation is conceptually misleading.

Although homogeneous solutions exist, their physical realization requires global conditions that cannot be generated by the electromagnetic field itself. Radiation is not produced dynamically by Maxwell’s equations but emerges only through boundary conditions, interaction with matter, and the global structure of spacetime.

This perspective unifies insights from constrained Hamiltonian dynamics \cite{Dirac1964,HenneauxTeitelboim}, relativistic interaction theorems \cite{Currie1963}, quantum field theory \cite{Haag1992}, absorber theory \cite{WheelerFeynman1945}, and gravitational radiation theory \cite{Bondi1962,AshtekarStreubel}.

\section{Maxwell Theory as a Constrained System}

The Maxwell action in Minkowski spacetime is
\begin{equation}
S = -\frac{1}{4} \int F_{\mu\nu} F^{\mu\nu} , d^4x + \int j^\mu A_\mu , d^4x .
\end{equation}

Canonical momenta are
\begin{equation}
\pi^i = F^{0i} = E^i, \qquad \pi^0 = 0,
\end{equation}
yielding a primary constraint. The Hamiltonian reads
\begin{equation}
H = \int d^3x \left[ \frac{1}{2}(\mathbf{E}^2 + \mathbf{B}^2) + A_0(\nabla \cdot \mathbf{E} - \rho) \right].
\end{equation}

Gauss’s law,
\begin{equation}
\nabla \cdot \mathbf{E} = \rho,
\end{equation}
is therefore a constraint rather than a dynamical equation. Only the transverse components of the field constitute physical degrees of freedom.

Crucially, these degrees of freedom evolve autonomously but do not self-generate. If absent initially, they remain absent for all time.

\section{Propagation Without Generation}

The hyperbolic nature of Maxwell’s equations ensures propagation of initial data but provides no mechanism for its creation. Radiation must therefore be encoded in initial or boundary conditions rather than generated dynamically.

This distinction is often obscured by treating homogeneous solutions as physical objects rather than as mathematical continuations of prescribed data.

\section{Gauge Redundancy and Physical Content}

Gauge invariance identifies entire classes of field configurations as physically equivalent. Only gauge-invariant quantities, such as the field strength, represent observables.

Since gauge symmetry eliminates longitudinal and scalar modes, electromagnetic radiation consists solely of transverse degrees of freedom that possess no self-interaction. This sharply limits the autonomous dynamical content of the theory.

\section{The No-Interaction Theorem}

The Currie–Jordan–Sudarshan theorem \cite{Currie1963} establishes that no relativistic particle theory with canonical coordinates and Poincaré invariance can exhibit nontrivial interactions. Interaction therefore requires field degrees of freedom.

However, electromagnetic fields themselves do not self-interact. Interaction is thus displaced into a relational structure between matter and field rather than residing in either alone.

\section{Haag’s Theorem and the Status of Free Fields}

Haag’s theorem \cite{Haag1992} demonstrates that interacting and free quantum fields cannot be related by a unitary transformation. Consequently, the interaction picture does not exist globally.

Free photons therefore do not exist as physical states within interacting quantum electrodynamics. They appear only as asymptotic constructs used to define scattering amplitudes.

\section{Asymptotic Structure and Boundary Conditions}

Radiation is defined operationally through asymptotic conditions at null infinity. The Sommerfeld radiation condition,
\begin{equation}
F_{\text{in}} = 0 \quad \text{on } \mathcal{I}^-,
\end{equation}
is not a consequence of local field equations but a global assumption about spacetime.

This choice introduces time asymmetry and encodes the presence of absorbers in the future.

\section{Radiation and the Arrow of Time}

Maxwell’s equations are time-reversal invariant, yet radiation is not. The asymmetry arises from boundary conditions reflecting the low-entropy past of the universe. Radiation is therefore inseparable from thermodynamic irreversibility.

\section{Wheeler--Feynman Absorber Theory}

In the absorber theory of Wheeler and Feynman \cite{WheelerFeynman1945}, interactions are fundamentally time-symmetric. Radiation reaction emerges only when the universe acts as a perfect absorber.

This framework makes explicit that radiation is not a local phenomenon but a global one, depending on the structure of the entire universe.

\section{Gravitational Radiation}

An analogous situation arises in general relativity. Gravitational radiation is defined only in asymptotically flat spacetimes via Bondi news \cite{Bondi1962,AshtekarStreubel}. In spatially closed universes, gravitational radiation is undefined.

Thus, radiation in both electromagnetism and gravity is intrinsically tied to asymptotic structure.

\section{Conclusion}

Radiation is not an autonomous physical entity. It is not generated locally, nor does it exist independently of global structure. Instead, it is an emergent, relational phenomenon arising from constraints, boundary conditions, and the impossibility of isolating subsystems in relativistic physics.

The homogeneous Maxwell equations permit radiation mathematically, but physics permits it only under specific global conditions. Radiation is therefore not something that happens \emph{in} spacetime; it is something spacetime allows.

\begin{thebibliography}{99}

\bibitem{Dirac1964}
P. A. M. Dirac, \textit{Lectures on Quantum Mechanics}, Yeshiva University Press (1964).

\bibitem{HenneauxTeitelboim}
M. Henneaux and C. Teitelboim, \textit{Quantization of Gauge Systems}, Princeton University Press (1992).

\bibitem{Currie1963}
D. G. Currie, T. F. Jordan, and E. C. G. Sudarshan, ``Relativistic invariance and Hamiltonian theories of interacting particles,'' \textit{Rev. Mod. Phys.} \textbf{35}, 350 (1963).

\bibitem{Haag1992}
R. Haag, \textit{Local Quantum Physics}, 2nd ed., Springer (1992).

\bibitem{WheelerFeynman1945}
J. A. Wheeler and R. P. Feynman, ``Interaction with the Absorber as the Mechanism of Radiation,'' \textit{Rev. Mod. Phys.} \textbf{17}, 157 (1945).

\bibitem{Bondi1962}
H. Bondi, M. G. J. van der Burg, and A. W. K. Metzner, ``Gravitational Waves in General Relativity VII. Waves from Axi-Symmetric Isolated Systems,'' \textit{Proc. R. Soc. Lond. A} \textbf{269}, 21 (1962).

\bibitem{AshtekarStreubel}
A. Ashtekar and M. Streubel, ``Symplectic geometry of radiative modes and conserved quantities at null infinity,'' \textit{Proc. R. Soc. Lond. A} \textbf{376}, 585 (1981).

\end{thebibliography}

\end{document}
